\begin{explanationbox}
	\textbf{\textcolor{CExplanation}{一句话直觉}}

	将 $\triangle ABC$ 沿内心 $I$ 分割为三个小三角形，每个小三角形的高均为 $r$，底边之和为周长 $a+b+c$，故总面积 $S = \frac12(a+b+c)r$，即 $S = pr$。

	\medskip
	\textbf{\textcolor{CExplanation}{核心拆解}}
	\begin{description}[style=nextline, leftmargin=2.8em, labelsep=0.5em, itemsep=0.45em]
		\item[\textbf{要点一（面积分割）：}] 连接内心 $I$ 与三顶点 $A,B,C$，得 $\triangle IBC,\triangle ICA,\triangle IAB$。
		\item[\textbf{要点二（统一高）：}] 因 $I$ 到各边距离均为 $r$，故各小三角形以原边为底的高都是 $r$。
		\item[\textbf{要点三（求和化简）：}] 三面积相加：$\frac12 a r+\frac12 b r+\frac12 c r = \frac12(a+b+c)r$，即 $pr$。
	\end{description}

	\medskip
	\textbf{\textcolor{CExplanation}{几何本质（等距分割）}}

	$I$ 是内心，到三边等距，使得分割的小三角形具有相同的高 $r$，面积之和即为等高的底边之和的一半。

	\medskip
	\textbf{\textcolor{CExplanation}{代数意义（线性关系）}}

	$S$ 与 $p$、$r$ 的关系是线性齐次的，即 $S = p\cdot r$，可看作 $S$ 关于 $p$ 或 $r$ 的正比例函数（另一量固定时）。

	\medskip
	\textbf{\textcolor{CExplanation}{考点价值（考法归纳）}}
	\begin{itemize}[leftmargin=*, itemsep=0.5em, topsep=0.4em]
		\item \textbf{考法一（求内切圆半径）：} 已知三边，先用海伦公式求面积 $S$，再代 $r = S/p$。
		\item \textbf{考法二（求三角形面积）：} 已知内切圆半径 $r$ 及半周长 $p$，直接得 $S=pr$。
	\end{itemize}

	\medskip
	\textbf{\textcolor{CExplanation}{顿悟点}}

	\textbf{$\displaystyle S=pr$ 将三角形面积转化为以半周长 $p$ 为长、$r$ 为宽的矩形面积，几何直观是三角形与矩形等积。}

	\medskip
	\textbf{\textcolor{CExplanation}{使用场景}}
	\begin{itemize}[leftmargin=*, itemsep=0.5em, topsep=0.4em]
		\item \textbf{场景一（直接套用）：} 给出三边与内切圆半径相关量，求面积或半径。
		\item \textbf{场景二（方程桥梁）：} 在解三角形综合题中，联系面积公式 $S=\frac12ab\sin C$ 与 $S=pr$，建立方程求边长或角度。
	\end{itemize}
\end{explanationbox}
